\section{The complete paid recursion}
\label{sec:recursion}

\subsection{Every child-table entry has a cost}

For a group count $r$, a singleton factor is incident with $2(r-1)$
directed pair lists. Applying the common entropy inequality to their
indicator gives the entry-count coefficient
\[
 \alpha_r=\frac{\sqrt{2(r-1)}}r.
\]
Filling a singleton table invokes one recursive child contraction for
every entry. A locally certified aggregation coefficient $\beta_r$
therefore gives the paid recurrence
\begin{equation}\label{eq:paid-recurrence}
 T_d\le\poly(s)\left[
       2^{(\alpha_{r_d}+o(1))K_d}T_{d+1}
       +2^{(\beta_{r_d}+o(1))K_d}\right],
 \qquad K_d=\kappa\frac{M_d}{\sqrt{u_d}}.
\end{equation}
Using only the smaller local aggregation term would omit the main
recursive work. We retain the complete recurrence.

The root uses the joint certificate of Lemma~\ref{lem:joint}.
At depth one we retain the incoming thirteen-group three-plan certificate
$\beta_{13}^2=1572/4225$. In the tail we use the seven-group plan
$(0,1,2)\mid(3,4)\mid(5,6)$ with $\beta_7^2=632/1225$.
Those two local estimates use the previously reviewed $12/5$ matrix
allowance and have the same entry and storage conventions
\cite{trees,pairwise}. Their support checks are included in our replay.

\subsection{An explicit integer schedule}

Write $\lambda=\lceil\log_2s\rceil$, $B=\lceil6\log_2\lambda\rceil$.
Keep the preceding fastest packet's schedule:
\begin{align}\label{eq:schedule}
 r_0&=10,& \eta_0&=\frac{301}{20000},\notag\\
 r_1&=13,& \eta_1&=\frac{119}{10000},\notag\\
 r_d&=7\quad(d\ge2),&
 \eta_d&=\frac{1549}{324000}\left(\frac49\right)^{d-2}.
\end{align}
Set
\[
 u_d=\lfloor\eta_d\lambda\rfloor,\quad
 a_d=\sum_{j=0}^d(u_j+B),\quad L_d=2^{a_d},\quad Q_d=r_d^{a_d}.
\]
The cumulative retained fractions are
\[
 \eta_0+\eta_1=\frac{539}{20000},\qquad
 \sum_{d\ge0}\eta_d=\frac8{225}.
\]
Use padding with relative slack $\lambda^{-3}$ and stop at
$M\le H=\lceil s/\lambda^{3/4}\rceil$. For a finite exceptional range
of $s$, use direct evaluation. This is an eventual asymptotic certificate;
we do not assert that its small positive gate margins yield a useful
finite-size crossover.

\subsection{Degree, packing, and depth gates}

With side degrees bounded by $R,C$, the sufficient inherited packing gates
are
\begin{equation}\label{eq:packing-gates}
 QR/M=o(1),\qquad Q^2C(\log s)^2/M=o(1),\qquad LQ^2/M=o(1).
\end{equation}
Initially $R,C\le s^{9/10+o(1)}$. Subsequently the retained side stays
under that bound and the fresh side has degree at most $2^{a_{d-1}+1}$.

The exact integer inequalities
\[
 10^{1000}<2^{3322},\qquad 13^{1000}<2^{3701},\qquad 7^{100}<2^{281}
\]
give strict rational upper bounds on the relevant logarithms. The largest
old-degree exponents are consequently separated from one by
\begin{center}
\begin{tabular}{lll}
\toprule Stage & Exponent upper bound & Positive gap\\
\midrule
Root & $9/10+2\eta_0(3322/1000)$ & $39/5000000$\\
Depth one & $9/10+(539/20000)(3701/1000)$ & $5161/20000000$\\
Tail & $9/10+(8/225)(281/100)$ & $1/11250$\\
\bottomrule
\end{tabular}
\end{center}
The fresh-side and packet-capacity exponents are at most
\[
 \frac8{225}\left(1+2\cdot\frac{3701}{1000}\right)<1.
\]
All gaps are fixed positive constants. They absorb the additive
$O((\log\lambda)^2)$ guards, all fixed polylogarithmic factors, and
$s/M\le\lambda^{3/4}$ on active calls. Fresh degrees remain below
$s^{8/225+o(1)}<s^{9/10}$.

For a noncircular depth argument, provisionally consider depths through
$\lceil(3/4)\log_3\lambda\rceil+O(1)$. Since $4/9>1/3$, the smallest
tail taper in that range is at least $\lambda^{1/4+o(1)}$.
The gates just proved therefore hold uniformly. Nearly lossless padding
then makes each tail child smaller than $M/3$ for all sufficiently large
$s$. An active path cannot exceed the provisional depth. This closes the
induction, proves all active tapers diverge, and makes their floors
negligible. The padding errors over $O(\log\lambda)$ levels multiply to
$1+o(1)$, yielding
\begin{equation}\label{eq:masses}
 M_0\le s,\quad M_1\le(1+o(1))s/5,\quad
 M_d\le(1+o(1))\frac{2s}{65}(2/7)^{d-2}\quad(d\ge2).
\end{equation}

\subsection{Precision, entry access, and leaves}

For a retained valuation $i<u_d$,
$D=2^{a_d-i}\ge2^{a_{d-1}+B+1}$. The inherited creation-time precision
bound is $p=O(s\lambda^2 2^{a_{d-1}})$. Hence
\[
 \frac{p}{MD/u_d}=O(\lambda^{15/4}2^{-B})=o(1).
\]
The additional Hermite terms are absorbed by the same buffer. The local
precision bound and entropy estimate thus apply uniformly, while
$a_d\le(8/225)\lambda+O((\log\lambda)^2)$ keeps every ring rank and
precision polynomial in $s$.

Original entry routines have cost $2^{O(s^{9/10})}$ times a polynomial;
fresh packet entry routines have cost $\exp(s^{8/225+o(1)})$. Direct leaves
cost $2^{O(s/\lambda^{3/4})}$ times these entry costs. Each is lower
order on the scale $s/\sqrt\lambda$. The immutable entry descriptions
prevent expansion of an entire ancestor contraction inside an entry call.
We first count arithmetic and entry requests, then charge their uniform
description-access cost; the input-length polynomial is paid once.
Polynomial overhead in $s$ over $O(\log\lambda)$ levels adds
$O(\lambda\log\lambda)=o(s/\sqrt\lambda)$ to the exponent.

\subsection{Unrolling every time branch and the simultaneous space bound}

On the normalized scale define
\[
 K_0^*=\frac\kappa{\sqrt{\eta_0}},\qquad
 K_1^*=\frac\kappa{5\sqrt{\eta_1}},\qquad
 K_2^*=\frac{2\kappa}{65\sqrt{\eta_2}}.
\]
The tail ratio is $\gamma=(2/7)/\sqrt{4/9}=3/7$. Its entry sum has
coefficient $\alpha_7/(1-\gamma)=\sqrt3/2>\beta_7$. A tail aggregation
at offset $t$ is bounded by a convex combination of the initial tail
aggregation and the infinite entry sum. Therefore the only candidates
needed to bound every branch of \eqref{eq:paid-recurrence} are
\begin{align}
 C_0&=\beta_{10}K_0^*
 =\frac\pi{\ln2}\sqrt{\frac{2\beta_{10}^2}{301/20000}}
 =\NewTime\ldots,\label{eq:C0}\\
 C_1&=\alpha_{10}K_0^*+\beta_{13}K_1^*
 =\frac\pi{\ln2}\left(\sqrt{\frac{7200}{301}}
                       +\sqrt{\frac{50304}{20111}}\right)
 =29.3352039993294\ldots,\label{eq:C1}\\
 C_\infty&=\alpha_{10}K_0^*+\alpha_{13}K_1^*
                 +\frac{\alpha_7K_2^*}{1-\gamma}\notag\\
 &=\frac\pi{\ln2}\left(\sqrt{\frac{7200}{301}}
        +\sqrt{\frac{19200}{20111}}
        +\sqrt{\frac{77760}{261781}}\right)
 =29.065748524641\ldots .\label{eq:Cinf}
\end{align}
Rational interval arithmetic verifies $C_0>C_1,C_\infty$ and
$C_0<29.747$. At the root the joint selector bounds local storage by
$(20/21)C_0$. At subsequent depths the local storage coefficients are at
most $7.16817894875$ and $2.04877062910$, followed by a decreasing
geometric tail. Caches on one active recursion path coexist, so their
\emph{sizes add}: their leading exponent is the maximum, not the sum.

\begin{theorem}[Improved normalized contraction bound]\label{thm:global}
Under the input promises and imported exact interfaces of
Section~\ref{sec:interface}, the described deterministic contraction
procedure has simultaneous bounds
\begin{align*}
 T(s,\mathcal L)&\le(s+\mathcal L)^{O(1)}
    2^{(C_T+o(1))s/\sqrt{\lceil\log_2s\rceil}},\\
 S(s,\mathcal L)&\le(s+\mathcal L)^{O(1)}
    2^{(C_S+o(1))s/\sqrt{\lceil\log_2s\rceil}},
\end{align*}
where
\[
 \boxed{C_T=\NewTime\ldots,\qquad
 C_S=\frac{20}{21}C_T=\NewSpace\ldots.}
\]
The safe eventual coefficients $29.747$ and $28.330$ are valid.
\end{theorem}
\begin{proof}
The exact plan changes only the association and order of the finite sums
and products in \eqref{eq:finite-sum}. The imported identity and recovery
therefore give the original contraction. The gate and depth argument
justifies every call. The finite certificates and the paid recurrence give
the three time candidates, whose maximum is $C_0$. The simultaneous
storage argument gives $C_S$. All remaining costs and uniform errors have
been charged above. Strict gaps to the displayed safe coefficients absorb
the lower-order terms for sufficiently large $s$.
\end{proof}
